\documentclass[11pt]{article}
\usepackage[margin=1in]{geometry}
\usepackage{amsmath,amssymb,amsthm,hyperref}
\newtheorem{theorem}{Theorem}
\newtheorem{lemma}{Lemma}
\newcommand{\dd}{\,d}
\newcommand{\R}{\mathbb R}
\title{All-order monotonicity of iterated parabolic-cylinder ratios\\
\large MathDB 364320: a spectral proof}
\author{Principia Math In Progress}
\date{17 September 2026; source audit updated 22 September 2026}
\begin{document}
\maketitle

\paragraph{Task 1: the original problem.}
For $n>1/2$ and $x\in\R$, set
\[
 R_n^{[1]}(x)=\frac{U(n-1,x)}{U(n,x)},\qquad
 R_n^{[k+1]}(x)=\frac{R_n^{[k]}(x)}{R_{n+1}^{[k]}(x)}.
\]
Here $U$ is the standard parabolic cylinder function. We prove the substantive
all-order assertions of Segura's Conjecture 3.6 in \cite{Segura}.
The comparison between successive orders starts at $k=2$:
the unqualified comparison printed in the source cannot include $k=1$,
since $R_n^{[1]}(x)\sim x$ and $R_n^{[2]}(x)\to1$ as $x\to+\infty$.
We do not use that indexing issue as a counterexample or as research credit.

\begin{theorem}\label{thm:main}
For every real $n>1/2$, every real $x$, and every integer $k\ge1$,
$R_n^{[k]}(x)>0$ and
\[
 \frac{d}{dx}R_n^{[k]}(x)>0.
\]
For every integer $k\ge2$,
\[
 0<R_n^{[k]}(x)<R_n^{[k+1]}(x)<1,
 \qquad \lim_{x\to+\infty}R_n^{[k]}(x)=1.
\]
Furthermore, for every fixed $n>1/2$ and finite real $x$,
\[
 \lim_{k\to\infty}R_n^{[k]}(x)=1.
\]
\end{theorem}

\paragraph{Publication scope.}
The proof below is complete for these quantifiers. All 23 pages of the
author's manuscript announced as forthcoming in AMS Contemporary Mathematics
have now been read; Conjecture 3.6 remains unresolved there. Acceptance is
inferred from the explicit author announcement, supported by the reviewer
acknowledgement and publication list. The audit assigns one new campaign
credit on that source basis. The separate typeset publisher chapter remains
unread; no identity with it is asserted. The accompanying audit records this
revision of the earlier stricter publisher-copy requirement and the bounded
public-claim checks. The mathematical proof is unchanged.

\section*{A positive spectral representation}
Write $D_{-a}(t)=U(a-1/2,t)$. For $a>0$,
\begin{equation}\label{eq:integral}
 D_{-a}(t)=\frac{e^{-t^2/4}}{\Gamma(a)}
 \int_0^\infty s^{a-1}e^{-s^2/2-ts}\dd s>0,
 \qquad D_0(t)=e^{-t^2/4}.
\end{equation}
In particular all ratios in the theorem are defined and positive. The
integral also gives Gaussian decay, with at most polynomial factors,
as $t\to+\infty$, including the first two $t$-derivatives. The standard
equation and derivative identity are
\begin{equation}\label{eq:ode}
 D_{-a}''(t)=\left(\frac{t^2}{4}-\frac12+a\right)D_{-a}(t),
 \qquad
 D_\nu'(t)=\frac t2D_\nu(t)-D_{\nu+1}(t).
\end{equation}
They can also be obtained by differentiating \eqref{eq:integral} and
integrating by parts; see \cite{DLMF} for the normalizations.

Fix a real endpoint $x$. Let $H_x$ be the Dirichlet self-adjoint realization
on $L^2(x,\infty)$ of
\[
 -\frac{d^2}{dt^2}+\frac{t^2}{4}-\frac12.
\]
More explicitly it is the Friedrichs operator of the closure, initially on
$C_c^\infty(x,\infty)$, of the quadratic form
\begin{equation}\label{eq:factor}
 \int_x^\infty\left(|u'|^2+
       \left(\frac{t^2}{4}-\frac12\right)|u|^2\right)\dd t
 =\int_x^\infty\left|u'+\frac t2u\right|^2\dd t.
\end{equation}
Thus $H_x\ge0$. It has no zero eigenvector: equality in
\eqref{eq:factor} would imply $u(t)=c e^{-t^2/4}$, and its Dirichlet
boundary value forces $c=0$.

Put
\[
 v_a(t)=\frac{D_{-a}(t)}{D_{-a}(x)},\qquad
 v_0(t)=e^{-(t^2-x^2)/4},\qquad
 m_x(a)=-\frac{D_{-a}'(x)}{D_{-a}(x)}.
\]
The function $v_0$ is a nonzero vector of $L^2(x,\infty)$; it is not
asserted to satisfy the Dirichlet boundary condition. Let $E_x$ be the
spectral resolution of $H_x$ and define the finite positive measure
\[
 \nu_x(B)=\langle v_0,E_x(B)v_0\rangle.
\]
It has positive total mass and is supported on $(0,\infty)$, since
$H_x\ge0$ and $\ker H_x=\{0\}$.

\begin{lemma}\label{lem:weyl}
For every $a>0$,
\begin{equation}\label{eq:mrep}
 m_x(a)=\frac x2+
 \int_{(0,\infty)}\frac{a\lambda}{\lambda+a}\dd\nu_x(\lambda).
\end{equation}
Consequently,
\begin{equation}\label{eq:increment}
 m_x(a+1)-m_x(a)=
 \int_{(0,\infty)}
 \frac{\lambda^2}{(\lambda+a)(\lambda+a+1)}\dd\nu_x(\lambda).
\end{equation}
\end{lemma}
\begin{proof}
Both $v_a$ and $v_0$ have boundary value $1$. Their difference has zero
boundary value, is square integrable, and, by \eqref{eq:ode} and the
decay stated above, belongs to the Dirichlet operator domain. The equation
for this difference is
\[
 (H_x+a)(v_a-v_0)=-a v_0.
\]
The resolvent exists for $a>0$, so
\begin{equation}\label{eq:resolvent}
 v_a=v_0-a(H_x+a)^{-1}v_0.
\end{equation}
On the other hand, the Wronskian
$W(t)=v_0(t)v_a'(t)-v_0'(t)v_a(t)$ satisfies
$W'(t)=a v_0(t)v_a(t)$. Its limit at infinity is zero, and
$W(x)=-m_x(a)+x/2$. Integration therefore gives
\[
 m_x(a)-\frac x2
 =a\langle v_0,v_a\rangle
 =a\langle v_0,[I-a(H_x+a)^{-1}]v_0\rangle.
\]
The spectral theorem and \eqref{eq:resolvent} give \eqref{eq:mrep}.
Its integral is finite because $0\le a\lambda/(\lambda+a)\le a$ and
$\nu_x$ is finite. Subtracting the two representations proves
\eqref{eq:increment}; its integrand lies between zero and one.
\end{proof}

\section*{An explicit positive formula at every order}
Set $a=n-1/2>0$. Since $R_n^{[1]}=D_{-a}/D_{-a-1}$,
\eqref{eq:increment} gives
\begin{equation}\label{eq:base}
 \frac{d}{dx}\log R_n^{[1]}(x)
 =m_x(a+1)-m_x(a)
 =\int\frac{\lambda^2}{(\lambda+a)(\lambda+a+1)}\dd\nu_x(\lambda).
\end{equation}
The measure is held fixed while taking differences in $a$. No derivative
of the endpoint-dependent measure with respect to $x$ is being taken.

Let $\Delta f(a)=f(a+1)-f(a)$ and $(b)_j=b(b+1)\cdots(b+j-1)$.
The defining ratio recurrence implies
\[
 \log R_n^{[k]}(x)=(-1)^{k-1}\Delta^{k-1}
                        \log R_n^{[1]}(x).
\]
For $b>0$ and $p\ge0$, the elementary identity
\begin{equation}\label{eq:fd}
 (-1)^p\Delta^p\frac1{b(b+1)}=\frac{(p+1)!}{(b)_{p+2}}
\end{equation}
follows by induction from
$1/(b)_j-1/(b+1)_j=j/(b)_{j+1}$.
Taking this finite linear combination in \eqref{eq:base} yields
\begin{equation}\label{eq:positive}
 \boxed{\quad
 \frac{d}{dx}\log R_n^{[k]}(x)=
 k!\int_{(0,\infty)}
 \frac{\lambda^2}{(\lambda+a)_{k+1}}\dd\nu_x(\lambda)>0.
 \quad}
\end{equation}
All integrals converge: the kernel for $k=1$ is bounded by one, and the
ratio of the kernel for $k+1$ to that for $k$ is
\begin{equation}\label{eq:kernelratio}
 \frac{k+1}{\lambda+a+k+1}<1.
\end{equation}
Strict positivity follows from the nonzero positive measure on $(0,\infty)$.
Since the ratios themselves are positive, \eqref{eq:positive} proves their
strict increase in $x$ for every $k\ge1$.

\section*{The upper bound and both limits}
In \eqref{eq:integral}, substitute $s=u/x$ for $x>0$.
Dominated convergence, with majorant $u^{a-1}e^{-u}$, proves
\[
 D_{-a}(x)\sim e^{-x^2/4}x^{-a},\qquad
 R_n^{[1]}(x)\sim x\qquad(x\to+\infty).
\]
For every fixed $k\ge2$, the expression for $\log R_n^{[k]}$ is a finite
difference of $\log x+o(1)$, so its limit is zero. Strict increase in $x$
then gives $\log R_n^{[k]}(x)<0$ for every finite $x$, proving
$R_n^{[k]}(x)<1$. The recurrence immediately gives
\[
 R_n^{[k+1]}(x)=R_n^{[k]}(x)/R_{n+1}^{[k]}(x)>R_n^{[k]}(x)
 \qquad(k\ge2).
\]

For completeness, the limit in the order also follows from the same
kernel. For fixed $b=\lambda+a>0$,
$k!/(b)_{k+1}\to0$: by \eqref{eq:kernelratio}, its successive ratios
are $(k+1)/(b+k+1)$, whose product tends to zero because
$\sum_k\log(1+b/(k+1))=\infty$.
Dominated convergence in \eqref{eq:positive} shows, for each fixed real $t$,
\[
 q_k(t):=\frac{d}{dt}\log R_n^{[k]}(t)\longrightarrow0,
 \qquad 0<q_k(t)\le q_2(t)\quad(k\ge2).
\]
Moreover, the already proved limit at $+\infty$ gives
\[
 -\log R_n^{[k]}(x)=\int_x^\infty q_k(t)\dd t,
 \qquad \int_x^\infty q_2(t)\dd t=-\log R_n^{[2]}(x)<\infty.
\]
A second application of dominated convergence proves
$\log R_n^{[k]}(x)\to0$. This completes the proof of Theorem
\ref{thm:main}.

\paragraph{Verification scope.}
This is an analytic proof for every required order and parameter.
The accompanying symbolic checker verifies the rational kernel identities
used in the proof; numerical checks verify conventions and offsets only.
Neither finite computation is used to infer an infinite sign assertion.
The audit was performed by the authoring agent, with no independent peer
review claimed.

\section*{Four follow-up tasks}
\paragraph{Task 2.}
Determine the sharp asymptotic size of $1-R_n^{[k]}(x)$ as $k\to\infty$
in terms of the lowest Dirichlet eigenvalue of $H_x$, including uniform
estimates as $n\downarrow1/2$ and $x\to-\infty$.
\paragraph{Task 3.}
Compute explicit spectral weights for $\nu_x$, including closed formulas
at $x=0$, and turn truncated versions of \eqref{eq:positive} into certified
two-sided bounds with effective tail estimates.
\paragraph{Task 4.}
Identify other contiguous special-function ratios admitting the same
positive-resolvent argument, specifying the parameter ranges in which the
associated Dirichlet operator is nonnegative and the required ratio
identities hold.
\paragraph{Task 5.}
Determine whether each $R_n^{[k]}(x)$ with $k\ge2$ has exactly one
inflection point, and how that point depends on $n$ and $k$.

\begin{thebibliography}{9}
\bibitem{Segura} J. Segura, \emph{On bounds for ratios of contiguous
hypergeometric functions}, arXiv:2408.05573v1 (2024), Conjecture 3.6,
p.~9. All 23 manuscript pages were reviewed. Published in
\emph{Contemporary Mathematics} \textbf{818} (2025),
71--93, DOI \href{https://doi.org/10.1090/conm/818/16368}
{10.1090/conm/818/16368}. Published Conjecture 3.6 is on p.~79;
only searchable excerpts of the published chapter were accessible.
\bibitem{DLMF} NIST Digital Library of Mathematical Functions,
\href{https://dlmf.nist.gov/12.5}{\S12.5} and
\href{https://dlmf.nist.gov/12.8}{\S12.8}, integral, recurrence, and
derivative conventions for parabolic cylinder functions.
\end{thebibliography}
\end{document}
